\documentclass{amsart}
% For dealing with references we use the comment environment
\usepackage{verbatim}
\newenvironment{reference}{\comment}{\endcomment}
%\newenvironment{reference}{}{}
\newenvironment{slogan}{\comment}{\endcomment}
\newenvironment{history}{\comment}{\endcomment}

% For commutative diagrams we use Xy-pic
\usepackage[all]{xy}

% We use 2cell for 2-commutative diagrams.
\xyoption{2cell}
\UseAllTwocells

% We use multicol for the list of chapters between chapters
\usepackage{multicol}

% This is generally recommended for better output
\usepackage{lmodern}
\usepackage[T1]{fontenc}

% For cross-file-references

% Package for hypertext links:
\usepackage{hyperref}

% For any local file, say "hello.tex" you want to link to please

% Theorem environments.
%
\theoremstyle{plain}
\newtheorem{theorem}[subsection]{Theorem}
\newtheorem{proposition}[subsection]{Proposition}
\newtheorem{lemma}[subsection]{Lemma}

\theoremstyle{definition}
\newtheorem{definition}[subsection]{Definition}
\newtheorem{example}[subsection]{Example}
\newtheorem{exercise}[subsection]{Exercise}
\newtheorem{situation}[subsection]{Situation}

\theoremstyle{remark}
\newtheorem{remark}[subsection]{Remark}
\newtheorem{remarks}[subsection]{Remarks}

\numberwithin{equation}{subsection}

% Macros
%
\def\lim{\mathop{\mathrm{lim}}\nolimits}
\def\colim{\mathop{\mathrm{colim}}\nolimits}
\def\Spec{\mathop{\mathrm{Spec}}}
\def\Hom{\mathop{\mathrm{Hom}}\nolimits}
\def\Ext{\mathop{\mathrm{Ext}}\nolimits}
\def\SheafHom{\mathop{\mathcal{H}\!\mathit{om}}\nolimits}
\def\SheafExt{\mathop{\mathcal{E}\!\mathit{xt}}\nolimits}
\def\Sch{\mathit{Sch}}
\def\Mor{\mathop{\mathrm{Mor}}\nolimits}
\def\Ob{\mathop{\mathrm{Ob}}\nolimits}
\def\Sh{\mathop{\mathit{Sh}}\nolimits}
\def\NL{\mathop{N\!L}\nolimits}
\def\CH{\mathop{\mathrm{CH}}\nolimits}
\def\proetale{{pro\text{-}\acute{e}tale}}
\def\etale{{\acute{e}tale}}
\def\QCoh{\mathit{QCoh}}
\def\Ker{\mathop{\mathrm{Ker}}}
\def\Im{\mathop{\mathrm{Im}}}
\def\Coker{\mathop{\mathrm{Coker}}}
\def\Coim{\mathop{\mathrm{Coim}}}

% Boxtimes
%
\DeclareMathSymbol{\boxtimes}{\mathbin}{AMSa}{"02}

%
% Macros for moduli stacks/spaces
%
\def\QCohstack{\mathcal{QC}\!\mathit{oh}}
\def\Cohstack{\mathcal{C}\!\mathit{oh}}
\def\Spacesstack{\mathcal{S}\!\mathit{paces}}
\def\Quotfunctor{\mathrm{Quot}}
\def\Hilbfunctor{\mathrm{Hilb}}
\def\Curvesstack{\mathcal{C}\!\mathit{urves}}
\def\Polarizedstack{\mathcal{P}\!\mathit{olarized}}
\def\Complexesstack{\mathcal{C}\!\mathit{omplexes}}
% \Pic is the operator that assigns to X its picard group, usage \Pic(X)
% \Picardstack_{X/B} denotes the Picard stack of X over B
% \Picardfunctor_{X/B} denotes the Picard functor of X over B
\def\Pic{\mathop{\mathrm{Pic}}\nolimits}
\def\Picardstack{\mathcal{P}\!\mathit{ic}}
\def\Picardfunctor{\mathrm{Pic}}
\def\Deformationcategory{\mathcal{D}\!\mathit{ef}}

\setlength{\emergencystretch}{3em}
\begin{document}
\title[Stacks Project, Tag 0BEM]{587 --- Euler characteristics of coherent sheaves twisted by several line bundles on proper schemes are numerical polynomials bounded by support dimension}
\maketitle
\noindent Copyright (C) 2005--2025 Johan de Jong. Stacks Project authors. Source: \href{https://stacks.math.columbia.edu/tag/0BEM}{Tag 0BEM}.

\noindent Editorial difficulty note (not author text). Difficulty H2: The proof must establish polynomiality without assuming projectivity or positivity of the line bundles. It removes embedded points, passes to the annihilator subscheme, and uses a regular meromorphic section with denominator ideal to express every finite difference as an Euler function on strictly smaller-dimensional support..

\noindent Editorial provenance note (not author text). History: excerpt from revision \texttt{a0444\allowbreak 6e57e\allowbreak c1fbc\allowbreak 252a8\allowbreak 71afc\allowbreak ec775\allowbreak 2fb28\allowbreak 07b14}, varieties.tex, lines 10194--10282. Reference presentation was adapted; original-author context and core support blocks were retained or added. The mathematical wording is unchanged. Licensed under GNU FDL 1.2 or later; no invariant sections or cover texts. See ../licenses/COPYING. Context blocks reproduce author definitions and supporting results; they are not additional counted problems.

\begin{lemma}
\label{lemma-numerical-polynomial-from-euler}
Let $k$ be a field. Let $X$ be a proper scheme over $k$. Let $\mathcal{F}$
be a coherent $\mathcal{O}_X$-module. Let
$\mathcal{L}_1, \ldots, \mathcal{L}_r$ be invertible $\mathcal{O}_X$-modules.
The map
$$
(n_1, \ldots, n_r) \longmapsto
\chi(X, \mathcal{F} \otimes
\mathcal{L}_1^{\otimes n_1} \otimes \ldots \otimes
\mathcal{L}_r^{\otimes n_r})
$$
is a numerical polynomial in $n_1, \ldots, n_r$ of total degree at
most the dimension of the support of $\mathcal{F}$.
\end{lemma}

\begin{proof}
We prove this by induction on $\dim(\text{Supp}(\mathcal{F}))$.
If this number is zero, then the function is constant with value
$\dim_k \Gamma(X, \mathcal{F})$ by Lemma \href{https://stacks.math.columbia.edu/tag/0AYT}{lemma chi tensor finite (Tag 0AYT)}.
Assume $\dim(\text{Supp}(\mathcal{F})) > 0$.

\medskip\noindent
If $\mathcal{F}$ has embedded associated points, then we can consider
the short exact sequence
$0 \to \mathcal{K} \to \mathcal{F} \to \mathcal{F}' \to 0$
constructed in Divisors, Lemma \href{https://stacks.math.columbia.edu/tag/02OL}{lemma remove embedded points (Tag 02OL)}.
Since the dimension of the support of $\mathcal{K}$ is strictly less,
the result holds for $\mathcal{K}$ by induction hypothesis and with
strictly smaller total degree.
By additivity of the Euler characteristic
(Lemma \href{https://stacks.math.columbia.edu/tag/08AA}{lemma euler characteristic additive (Tag 08AA)})
it suffices to prove the result for $\mathcal{F}'$. Thus we may assume
$\mathcal{F}$ does not have embedded associated points.

\medskip\noindent
If $i : Z \to X$ is a closed immersion and $\mathcal{F} = i_*\mathcal{G}$,
then we see that the result for $X$, $\mathcal{F}$,
$\mathcal{L}_1, \ldots, \mathcal{L}_r$ is equivalent to the result
for $Z$, $\mathcal{G}$, $i^*\mathcal{L}_1, \ldots, i^*\mathcal{L}_r$
(since the cohomologies agree, see
Cohomology of Schemes, Lemma \href{https://stacks.math.columbia.edu/tag/089W}{lemma relative affine cohomology (Tag 089W)}).
Applying Divisors, Lemma \href{https://stacks.math.columbia.edu/tag/02OM}{lemma no embedded points endos (Tag 02OM)}
we may assume that $X$ has no embedded components and
$X = \text{Supp}(\mathcal{F})$.

\medskip\noindent
Pick a regular meromorphic section $s$ of $\mathcal{L}_1$, see
Divisors, Lemma \href{https://stacks.math.columbia.edu/tag/02OZ}{lemma regular meromorphic section exists (Tag 02OZ)}.
Let $\mathcal{I} \subset \mathcal{O}_X$ be the ideal of
denominators of $s$ and consider the maps
$$
\mathcal{I}\mathcal{F} \to \mathcal{F},\quad
\mathcal{I}\mathcal{F} \to \mathcal{F} \otimes \mathcal{L}_1
$$
of Divisors, Lemma \href{https://stacks.math.columbia.edu/tag/02P2}{lemma make maps regular section (Tag 02P2)}.
These are injective and have cokernels $\mathcal{Q}$, $\mathcal{Q}'$
supported on nowhere dense closed subschemes of $X = \text{Supp}(\mathcal{F})$.
Tensoring with the invertible module
$\mathcal{L}_1^{\otimes n_1} \otimes \ldots \otimes \mathcal{L}_r^{\otimes n_r}$
is exact, hence using additivity again
we see that
\begin{align*}
&\chi(X, \mathcal{F} \otimes \mathcal{L}_1^{\otimes n_1} \otimes \ldots \otimes
\mathcal{L}_r^{\otimes n_r}) -
\chi(X, \mathcal{F} \otimes \mathcal{L}_1^{\otimes n_1 + 1}
\otimes \ldots \otimes \mathcal{L}_r^{\otimes n_r}) \\
& =
\chi(\mathcal{Q} \otimes \mathcal{L}_1^{\otimes n_1} \otimes \ldots \otimes
\mathcal{L}_r^{\otimes n_r}) -
\chi(\mathcal{Q}' \otimes \mathcal{L}_1^{\otimes n_1} \otimes \ldots \otimes
\mathcal{L}_r^{\otimes n_r})
\end{align*}
Thus we see that the function $P(n_1, \ldots, n_r)$ of the lemma has
the property that
$$
P(n_1 + 1, n_2, \ldots, n_r) - P(n_1, \ldots, n_r)
$$
is a numerical polynomial of total degree $<$ the dimension
of the support of $\mathcal{F}$. Of course by symmetry the same
thing is true for
$$
P(n_1, \ldots, n_{i - 1}, n_i + 1, n_{i + 1}, \ldots, n_r)
- P(n_1, \ldots, n_r)
$$
for any $i \in \{1, \ldots, r\}$. A simple arithmetic argument shows
that $P$ is a numerical polynomial of total degree at most
$\dim(\text{Supp}(\mathcal{F}))$.
\end{proof}
\end{document}
